\documentclass[11pt]{article}
\usepackage[margin=1.1in]{geometry}
\usepackage{amsmath,amssymb,amsthm}
\usepackage[colorlinks=true,linkcolor=blue,citecolor=blue]{hyperref}

\newtheorem{theorem}{Theorem}[section]
\newtheorem{lemma}[theorem]{Lemma}
\newtheorem{proposition}[theorem]{Proposition}
\newtheorem{corollary}[theorem]{Corollary}
\newtheorem{definition}[theorem]{Definition}
\newtheorem{remark}[theorem]{Remark}
\newtheorem{problem}[theorem]{Problem}

\newcommand{\Q}{Q_\infty}
\newcommand{\F}{\mathcal{F}}
\newcommand{\dist}{\operatorname{dist}}
\newcommand{\diam}{\operatorname{diam}}
\newcommand{\EX}{\mathsf{Ex}}
\newcommand{\W}{\mathcal{W}}
\newcommand{\Gcal}{\mathcal{G}}
\newcommand{\mdim}{\operatorname{mdim}_{\mathrm M}}

\title{Finite-horizon windowed distributional chaos\\
on the Hilbert cube: a residual theorem and its sharp\\
finite-certificate obstruction}
\author{Yitzchak Shmalo\\
\small Einstein Institute of Mathematics, The Hebrew University of Jerusalem, Jerusalem, Israel\\
\small }
\date{June 2026}

\begin{document}
\maketitle

\begin{abstract}
On the Hilbert cube $\Q=[0,1]^{\mathbb N}$ the companion paper \cite{companion}
establishes a residual depth hierarchy of margin-stable covering schemes in the full
non-autonomous mapping space $\F(\Q)=C(\Q,\Q)^{\mathbb N}$, together with an
empty-interior obstruction showing that no open class can certify exact covering in
arbitrarily deep coordinates; it leaves open whether the infinite-horizon dynamical
conclusions of the underlying dense theorem of \cite{Shmalo2026} --- all-tail
semifactors, windowed distributional chaos of type~1, infinite entropy and infinite
metric mean dimension --- are residual.  This paper advances that question on its
\emph{finite-horizon} layer.  First we prove a positive residual theorem: for any
countable family of dense-open window sequences there is a comeager set
$\Gcal\subseteq\F(\Q)$ such that every system in $\Gcal$ admits finite-horizon windowed
DC1 witnesses of every order, at every tail, every rational proximal scale, every
rational tolerance, and the fixed separation scale --- pairs of orbits proximal on an
arbitrarily large prescribed time-fraction of arbitrarily long horizons and separated on
a positive time-fraction, both visiting the prescribed windows up to those horizons.
The same residual set in fact contains finite scrambled packets of every prescribed
cardinality: for every \(M\) there are arbitrarily long horizons carrying \(M\) orbits
which are all mutually proximal on a common large set of times and whose every pair is
\(t_0\)-separated on a positive fraction of times.  The mechanism is a
finite-horizon stability lemma whose orbit-tracking modulus does \emph{not} decay with
coordinate depth, in contrast with the depth obstruction of \cite{companion}, combined
with the dense windowed-horseshoe theorem of \cite{Shmalo2026}, a block schedule, and a
binary coding argument across the separation blocks.
Second we prove that this is sharp at the level of certificates: every finite-horizon
certified open set --- one whose membership depends only on a bounded window of
coordinate maps --- that is contained in the class of systems with an
infinite-metric-mean-dimension tail, or in the corresponding scale-law entropy class from
\cite{Shmalo2026}, is \emph{empty}.  The proof is an explicit coordinate squeeze
applied past the certifying window, which stays uniformly close to every system in the
set yet caps the metric mean dimension of every later tail.  The same squeeze also defeats
the all-tail symbolic semifactor, while leaving a bare uncountable DC1 set untouched, so
the certificate obstruction separates the growth-type conclusions from the cardinality
conclusion.  We record the Banach--Mazur corollary and an honest ledger: this rules out the natural
finite-certificate method class for the infinite-horizon problem, but does not decide
whether the full conclusions are comeager; that remains open for methods that are not
finite-horizon.
\end{abstract}

\section{Introduction}\label{sec:intro}

A non-autonomous discrete dynamical system on a compact metric space $(X,\rho)$ is a
sequence $f=(f_1,f_2,\dots)$ of continuous self-maps acting by composition
$F^j=f_{j}\circ\cdots\circ f_1$.  The full mapping space
\[
\F(X)=C(X,X)^{\mathbb N},
\qquad
D_X(f,g)=\sup_{n\ge1}\,\sup_{x\in X}\rho\bigl(f_n(x),g_n(x)\bigr),
\]
is complete and hence Baire.  On finite-dimensional cubes and on compact manifolds the
construction of \cite{Shmalo2026} produces an open dense class of systems with all-tail
full complexity: after deleting any finite initial segment one still has a compact
symbolic semifactor onto a full variable-alphabet shift, a continuum-sized scrambled
set for distributional chaos of type~1 (DC1) visiting prescribed dense open windows at
prescribed times, infinite topological entropy, and the sharp small-scale entropy
spectrum.  On the Hilbert cube
\[
\Q=[0,1]^{\mathbb N},
\qquad
\rho(x,y)=\sum_{j\ge1}2^{-j}\,|x_j-y_j|,
\]
the same conclusions hold only \emph{densely}, because coordinate $j$ carries metric
weight $2^{-j}$ and exact control of arbitrarily deep coordinates is not stable.  The
companion paper \cite{companion} analyses exactly this gap.  It proves that for every
depth $N$ the class of systems carrying a margin-stable depth-$N$ covering scheme is
open dense, with an explicit stability radius of order $2^{-N}$, so that the
intersection over all depths is residual; that on this residual set one-step covering
is exact at every finite depth; and, on the negative side, that the dense
exact-horseshoe class --- indeed any class forcing some image to have full range in
infinitely many coordinates --- has empty interior.  What \cite{companion} explicitly
leaves open (its Problem, restated below as Problem~\ref{prob:companion}) is whether the
\emph{infinite-horizon dynamical} conclusions are themselves residual:

\begin{quotation}
\noindent``What is missing is exactly a mechanism producing, for a residual set of
systems, orbits threading infinitely many prescribed bricks --- equivalently, a
replacement for the backward induction \cite[Lemma~13.1]{Shmalo2026} that consumes only
finite-depth information per step.''
\end{quotation}

This open question traces back to Section~14 of \cite{Shmalo2026}, where the
``dense, not open'' phenomenon on $\Q$ is first recorded and a residual theorem ``by a
Baire-fusion mechanism'' is requested.

\subsection*{This paper}

We attack the question on the layer where it is tractable: \emph{finite horizons}.  The
infinite-horizon DC1 conclusion asks for a scrambled pair $(x,y)$ whose orbits are
proximal along a set of times of upper density~$1$ and separated along a set of upper
density~$1$, with prescribed window visits at \emph{all} times.  We weaken ``all
times'' to ``all times below a horizon $N$'' and ``upper density'' to ``a prescribed
fraction of $\{0,\dots,N-1\}$'', and prove that the resulting finite-horizon witnesses
exist residually, of every order and at every scale, and then prove that no
finite-horizon \emph{certificate} can reach the infinite-horizon conclusions.  The two
results sit on opposite sides of a sharp line.

\begin{theorem}[Residual finite-horizon windowed DC1 witnesses;
Theorem~\ref{thm:residual} below]\label{thm:A}
Let $\{W^a_\bullet\}_{a\ge1}$ be any countable family of sequences of dense open subsets
of $\Q$, and fix a separation scale \(t_0\in(0,1)\).  There is a comeager set
$\Gcal\subseteq\F(\Q)$ such that every $g\in\Gcal$ has the following property: for every
index $a$, every tail $r\ge0$, every rational proximal scale $\varepsilon\in(0,1)$, and
all rational $0<\theta<\eta<1$, for every $N_0\ge1$ there exists $N\ge N_0$ for which
there are points $x,y\in\Q$ and disjoint sets of times
$P,S\subseteq\{0,\dots,N-1\}$ with $\#P\ge(1-\eta)N$ and $\#S\ge\theta N$ for which
\[
\rho\bigl(G^j_r(x),G^j_r(y)\bigr)<\varepsilon\ \ (j\in P),
\qquad
\rho\bigl(G^j_r(x),G^j_r(y)\bigr)>t_0\ \ (j\in S),
\]
and $G^j_r(x),G^j_r(y)\in W^a_{r+j+1}$ for every $j<N$, where
$G^j_r=g_{r+j}\circ\cdots\circ g_{r+1}$ and $G^0_r=\mathrm{id}$.
\end{theorem}

The witness in Theorem~\ref{thm:A} is the finite-horizon, fixed-tolerance shadow of a
DC1 scrambled pair: proximal on a $(1-\eta)$-fraction of the horizon, separated at a
fixed scale on a positive $\theta$-fraction, inside the prescribed windows up to the
horizon.  Letting $\eta\downarrow0$ and then choosing any smaller $\theta>0$ produces
arbitrarily long horizons on which the pair is proximal on an arbitrarily large fraction
of times and still separated at scale $t_0$ on a positive fraction --- the finite-horizon
content of ``proximal and separated of every order.''

The same construction also gives finite scrambled packets, not only pairs.

\begin{theorem}[Residual finite-horizon scrambled packets]\label{thm:packet-intro}
Under the hypotheses of Theorem~\ref{thm:A}, possibly after replacing the comeager set
by a smaller comeager set, every \(g\) has the following stronger property.  For every
window sequence \(W^a_\bullet\), every tail \(r\ge0\), every integer \(M\ge2\), every
rational proximal scale \(\varepsilon\in(0,1)\), every rational \(0<\eta<1\), and every
\(N_0\), there are \(N\ge N_0\), points
\[
  x^1,\ldots,x^M\in\Q,
\]
a common set \(P\subseteq\{0,\ldots,N-1\}\), and pair-dependent sets
\(S_{uv}\subseteq\{0,\ldots,N-1\}\), \(1\le u<v\le M\), such that
\[
  \#P\ge(1-\eta)N,\qquad
  \#S_{uv}\ge \frac{\eta}{8}N\quad(1\le u<v\le M),
\]
all \(M\) orbit segments visit the prescribed windows up to time \(N\), all pairs are
\(\varepsilon\)-proximal at every time in \(P\), and the pair \((x^u,x^v)\) is
\(t_0\)-separated at every time in \(S_{uv}\).
\end{theorem}

\begin{theorem}[No finite-horizon certificate forces infinite-horizon complexity;
Theorem~\ref{thm:obstruction} below]\label{thm:B}
Call an open set $V\subseteq\F(\Q)$ \emph{finite-horizon certified} if there are an
offset $r\ge0$ and a horizon $N\ge1$ such that membership $g\in V$ depends only on the
finite window $(g_{r+1},\dots,g_{r+N})$.  If $V$ is finite-horizon certified and
\[
V\subseteq\bigl\{\,g\in\F(\Q):\ \mdim(\mathrm{tail}_s g)=\infty
\ \text{for some }s\ge0\,\bigr\},
\]
then $V=\varnothing$.  The same holds if the entropy hypothesis is replaced by the
scale-law entropy class of \cite{Shmalo2026}: for some tail and every \(q\), for some
constant \(C_q\), the
small-scale entropy satisfies
\[
h^{(\varepsilon)}_{\mathrm{top}}(\mathrm{tail}_s g)
\ge q\log(1/\varepsilon)-C_q
\]
for all sufficiently small \(\varepsilon\).
\end{theorem}

Theorem~\ref{thm:B} says the line drawn by Theorem~\ref{thm:A} cannot be crossed by the
obvious method.  A natural strategy for the infinite-horizon problem is to exhibit a
countable family of open dense sets, each certified by a bounded window, whose
intersection forces infinite-horizon complexity, and to invoke Baire.  Theorem~\ref{thm:B}
shows that every individual set in such a family that already entails an
infinite-\(\mdim\) tail, or the scale-law entropy tail supplied by the dense theorem, is
empty, hence not dense: the entailment and finite-horizon certification are
incompatible.  The infinite-horizon conclusions are genuinely
non-finite-horizon properties, and any residual proof must read unboundedly many
coordinate maps to certify a single system.

\subsection*{Why depth and horizon behave oppositely}

The contrast between the two theorems is the engine of the paper, and it is worth
isolating.  The obstruction of \cite{companion} is about \emph{coordinate depth}: a
condition that constrains image coordinate $\ell$ has $D$-stability only $O(2^{-\ell})$,
so margins erode geometrically in depth and no open set survives unbounded depth.  By
contrast, the finite-horizon witnesses of Theorem~\ref{thm:A} are stable because their
defining inequalities involve the metric $\rho$ \emph{directly}, not individual deep
coordinates, and the orbit-tracking modulus over a \emph{fixed} horizon $N$ is finite
and depth-free: if $D(h,g)<\delta$ then the two $N$-orbits stay within
$e_N\le\delta+\omega(e_{N-1})$, a quantity that depends on the uniform-continuity moduli
$\omega$ of the maps and on $N$, but not on how deep the proximity or separation is read
(Lemma~\ref{lem:fh-stability}).  This is the precise sense in which the finite-horizon
layer escapes the depth obstruction: horizons are stable while depths are not.  The two
obstructions are independent, and the present Theorem~\ref{thm:B} is a genuinely
\emph{horizon} obstruction --- the squeeze acts past a fixed window in time, not past a
fixed coordinate --- complementary to the depth obstruction of \cite{companion}.

\subsection*{An honest ledger}

We are careful about what is and is not proved.  Theorem~\ref{thm:A} produces witnesses
that are finite-horizon and carry a \emph{fixed} separation tolerance $t_0$; it does not
produce a single scrambled pair working at all horizons simultaneously, nor does it
recover the upper-density-one statements that define DC1 in the limit.  Theorem~\ref{thm:B}
rules out the natural finite-certificate route to the infinite-horizon problem, but it
does \emph{not} prove that the infinite-horizon class is non-comeager: a residual set
need not be a countable intersection of finite-horizon certified open sets.  Thus
Problem~\ref{prob:companion} (residuality of the full infinite-horizon conclusions)
remains open; what we add is a residual finite-horizon theorem on one side and, on the
other, a proof that the most direct method class cannot settle it.  Section~\ref{sec:open}
states this precisely.

\subsection*{Organization}

Section~\ref{sec:prelim} fixes notation, recalls the Baire property and the inputs we
import from \cite{Shmalo2026,companion}.  Section~\ref{sec:witness} defines the
finite-horizon windowed DC1 witness and proves its openness and depth-free stability.
It also defines finite scrambled packets and proves their openness.  Section~\ref{sec:density}
proves density via the dense windowed-horseshoe theorem, a block schedule, and a binary
coding argument across separation times, then assembles Theorems~\ref{thm:A} and
\ref{thm:packet-intro} by Baire intersection.
Section~\ref{sec:obstruction} proves the finite-horizon certificate obstruction
(Theorem~\ref{thm:B}) and its Banach--Mazur corollary.  Section~\ref{sec:open} records
the ledger and the open problem.

\section{Preliminaries}\label{sec:prelim}

Throughout, $\Q=[0,1]^{\mathbb N}$ carries
$\rho(x,y)=\sum_{j\ge1}2^{-j}|x_j-y_j|$, which induces the product topology and makes
$\Q$ compact, with $\diam_\rho(\Q)=1$.  We write $\pi_\ell:\Q\to[0,1]$ for the $\ell$-th
coordinate projection.  The space $C(\Q,\Q)$ carries the uniform metric
$d_\infty(u,v)=\sup_x\rho(u(x),v(x))$, and
\[
\F(\Q)=C(\Q,\Q)^{\mathbb N},\qquad
D(f,g)=\sup_{n\ge1}d_\infty(f_n,g_n).
\]
For $g\in\F(\Q)$, $r\ge0$ and $j\ge0$ put $G^0_r=\mathrm{id}$ and
$G^j_r=g_{r+j}\circ\cdots\circ g_{r+1}$; the \emph{tail at $r$} is
$\mathrm{tail}_rg=(g_{r+1},g_{r+2},\dots)$.

\begin{lemma}[{\cite[Lemma~2.1]{Shmalo2026}}]\label{lem:baire}
$(\F(\Q),D)$ is a complete metric space, and hence a Baire space.
\end{lemma}

\begin{proof}
$C(\Q,\Q)$ is complete because $\Q$ is compact and uniform limits of continuous maps
are continuous.  If $(f^{(k)})_k$ is $D$-Cauchy, each coordinate sequence
$(f^{(k)}_n)_k$ is $d_\infty$-Cauchy with a modulus independent of $n$, and the
coordinate limits assemble to $f\in\F(\Q)$ with $D(f^{(k)},f)\to0$.  Completeness gives
the Baire property.
\end{proof}

We use the elementary weight estimate of \cite[Eq.~(2.2)]{companion}: for
$u,v\in C(\Q,\Q)$ and any coordinate $\ell$,
\begin{equation}\label{eq:weight}
\sup_{x\in\Q}\bigl|\pi_\ell u(x)-\pi_\ell v(x)\bigr|\le2^{\ell}\,d_\infty(u,v),
\qquad
d_\infty(u,v)\le\sum_{\ell\ge1}2^{-\ell}\sup_x|\pi_\ell u(x)-\pi_\ell v(x)|.
\end{equation}

\paragraph{Imported dense theorem.}
The single dynamical input we need from \cite{Shmalo2026} is the exact windowed
horseshoe theorem on $\Q$, in the windowed-block form, together with the universal block
schedule.  We state it in the form used below.

\begin{theorem}[{Dense exact windowed horseshoe; \cite[Theorem~3.6 and \S12]{Shmalo2026}}]\label{thm:dense}
Let $W_\bullet=(W_n)_{n\ge1}$ be a sequence of dense open subsets of $\Q$.  For every
$f\in\F(\Q)$ and every $\eta_0>0$ there is $g\in\F(\Q)$ with $D(f,g)<\eta_0$ admitting
an exact $W_\bullet$-windowed Hilbert horseshoe: there are $\alpha\in(0,1/4)$, support
levels $K_n$, $(\alpha,K_n)$-boxes $J_n$, and finite-support boxes
$B_{n,1},\dots,B_{n,M_n}\subseteq\operatorname{int}(J_n)\cap W_n$ with
$\diam_\rho(B_{n,i})<\gamma_n\to0$ and $g_n(B_{n,i})=J_{n+1}$ for all $n,i$, organized on
a universal block schedule in which every label from the alphabet
$\{C_a,D_s,E_{q,m}\}$ occurs on infinitely many blocks.  On such a $g$, for every
$r\ge0$, the tail $\mathrm{tail}_rg$ carries a compact symbolic semifactor onto a full
variable-alphabet shift; in particular every finite admissible itinerary of bricks from
time $r+1$ onward is realized by an orbit of $\mathrm{tail}_rg$, with each visited brick
lying in its prescribed window.
\end{theorem}

The only facts we use from Theorem~\ref{thm:dense} are: (i) along an exact windowed
horseshoe, \emph{any} prescribed finite sequence of admissible bricks is followed
exactly by some orbit (the finite-itinerary realizability built into the symbolic
semifactor); and (ii) the schedule guarantees, on every long block of times, a
controllable proportion of ``proximity-type'' bricks $C_a$ that funnel two orbits into a
common small box and ``separation-type'' bricks $D_s$ that force them apart at the fixed
scale $t_0$.  We make (ii) quantitative in Section~\ref{sec:density}.

\paragraph{Companion inputs.}
From \cite{companion} we use only the framework (the metric $D$, the weight
estimate~\eqref{eq:weight}, the Baire property) and, in Section~\ref{sec:obstruction},
the empty-interior phenomenon as a point of contrast; the proofs below are otherwise
self-contained.

\section{The finite-horizon windowed DC1 witness}\label{sec:witness}

Fix a separation scale \(t_0\in(0,1)\) once and for all.

\begin{definition}[Finite-horizon windowed DC1 witness]\label{def:witness}
Let $W_\bullet=(W_n)$ be a sequence of dense open subsets of $\Q$, let $r\ge0$,
$\varepsilon\in(0,1)$, \(0<\theta<\eta<1\) and $N\ge1$.  We say $g\in\F(\Q)$ has the
\emph{finite-horizon windowed DC1 witness}
$\W_{r,\varepsilon,t_0,\eta,\theta,N}(W_\bullet)$ if there exist $x,y\in\Q$ and disjoint
$P,S\subseteq\{0,\dots,N-1\}$ with $\#P\ge(1-\eta)N$ and $\#S\ge\theta N$ such that, with
$G^j_r=g_{r+j}\circ\cdots\circ g_{r+1}$,
\begin{enumerate}
\item[(W1)] (\emph{proximity}) $\rho\bigl(G^j_r(x),G^j_r(y)\bigr)<\varepsilon$ for all
$j\in P$;
\item[(W2)] (\emph{separation}) $\rho\bigl(G^j_r(x),G^j_r(y)\bigr)>t_0$ for all
$j\in S$;
\item[(W3)] (\emph{windows}) $G^j_r(x)\in W_{r+j+1}$ and $G^j_r(y)\in W_{r+j+1}$ for all
$j$ with $0\le j<N$.
\end{enumerate}
Write $\W_{r,\varepsilon,t_0,\eta,\theta,N}(W_\bullet)\subseteq\F(\Q)$ for the set of systems
with this witness.
\end{definition}

The point of the next lemma is stability with a modulus that is finite for each $N$ and
\emph{does not decay with coordinate depth}.  This is the structural difference from the
depth obstruction of \cite[Remark~3.6]{companion}, where margins shrink like $2^{-N}$ in
the depth $N$; here $N$ is a \emph{time} horizon and the witness is open with a radius
that depends only on the horizon length and the uniform-continuity moduli of the
finitely many maps involved.

\begin{lemma}[Finite-horizon stability: openness]\label{lem:fh-stability}
Suppose $g$ has the witness \(\W_{r,\varepsilon,t_0,\eta,\theta,N}(W_\bullet)\) \emph{strictly},
meaning there is $\kappa>0$ such that {\rm(W1)} holds with $\varepsilon-2\kappa$ in
place of $\varepsilon$ on $P$, {\rm(W2)} holds with $t_0+2\kappa$ on $S$, and each
window point lies $\rho$-distance $>2\kappa$ inside its window $W_{r+j+1}$.  Then there
is $\delta>0$ such that every $h\in\F(\Q)$ with $D(h,g)<\delta$ also has the witness
\(\W_{r,\varepsilon,t_0,\eta,\theta,N}(W_\bullet)\), with the \emph{same} \(x,y,P,S\).  In
particular \(\W_{r,\varepsilon,t_0,\eta,\theta,N}(W_\bullet)\) is open in \((\F(\Q),D)\).
\end{lemma}

\begin{proof}
The maps $g_{r+1},\dots,g_{r+N-1}$ are uniformly continuous; let $\omega$ be a common
modulus of continuity (with respect to $\rho$) for all of them, so $\omega$ is
nondecreasing, $\omega(0^+)=0$, and $\rho(g_{r+i}(u),g_{r+i}(v))\le\omega(\rho(u,v))$
for $1\le i\le N-1$.  Fix $x$ (the same argument applies to $y$) and write
$X_j=G^j_r(x)$ for the $g$-orbit and $X^h_j=H^j_r(x)$ for the $h$-orbit, where
$H^j_r=h_{r+j}\circ\cdots\circ h_{r+1}$.  Set $e_j=\rho(X^h_j,X_j)$.  Then $e_0=0$ and,
for $1\le j<N$,
\[
e_j=\rho\bigl(h_{r+j}(X^h_{j-1}),g_{r+j}(X_{j-1})\bigr)
\le\rho\bigl(h_{r+j}(X^h_{j-1}),g_{r+j}(X^h_{j-1})\bigr)
+\rho\bigl(g_{r+j}(X^h_{j-1}),g_{r+j}(X_{j-1})\bigr)
\le D(h,g)+\omega(e_{j-1}).
\]
Writing $\delta_0=D(h,g)$, the recursion $e_j\le\delta_0+\omega(e_{j-1})$, $e_0=0$,
gives a bound $e_j\le E_j(\delta_0)$ where
$E_0=0$, $E_j(\delta_0)=\delta_0+\omega(E_{j-1}(\delta_0))$.  Each $E_j$ is a finite
composition of the continuous nondecreasing maps $s\mapsto\delta_0+\omega(s)$ evaluated
at the continuous-in-$\delta_0$ starting value, hence $E_j(\delta_0)\to0$ as
$\delta_0\to0$ for every fixed $j$.  Choose $\delta>0$ so small that
$E_{N-1}(\delta)<\kappa$; this is possible because there are only finitely many indices
$j\le N-1$ and each $E_j$ is continuous at $0$ with $E_j(0)=0$.  Then for $D(h,g)<\delta$
we have $\rho(X^h_j,X_j)<\kappa$ and, identically, $\rho(Y^h_j,Y_j)<\kappa$, for all
$j<N$.

It remains to transfer (W1)--(W3).  For $j\in P$,
\[
\rho(X^h_j,Y^h_j)\le\rho(X^h_j,X_j)+\rho(X_j,Y_j)+\rho(Y_j,Y^h_j)
<\kappa+(\varepsilon-2\kappa)+\kappa=\varepsilon,
\]
so (W1) holds for $h$.  For $j\in S$,
\[
\rho(X^h_j,Y^h_j)\ge\rho(X_j,Y_j)-\rho(X_j,X^h_j)-\rho(Y_j,Y^h_j)
>(t_0+2\kappa)-\kappa-\kappa=t_0,
\]
so (W2) holds.  For (W3), each $g$-orbit point lies $\rho$-distance $>2\kappa>\kappa$
inside its open window, and $\rho(X^h_j,X_j)<\kappa$, so $X^h_j\in W_{r+j+1}$, and
likewise for $y$.  Thus $h$ has the witness with the same data, proving openness; the
strict witness for $g$ is itself an instance, so the open set is nonempty when populated
as in Section~\ref{sec:density}.
\end{proof}

\begin{remark}[Depth-freeness, made explicit]\label{rem:depth-free}
The bound $e_j\le\delta_0+\omega(e_{j-1})$ never refers to a coordinate index: the only
inputs are the metric distances and the moduli $\omega$ of the finitely many maps
$g_{r+1},\dots,g_{r+N-1}$.  Consequently the stability radius $\delta$ of
Lemma~\ref{lem:fh-stability} is governed by the \emph{horizon} $N$ and the maps' moduli,
not by how deep a coordinate the proximity scale $\varepsilon$ or the separation scale
$t_0$ resolves.  Contrast \cite[Remark~3.6]{companion}: a covering condition reading
image coordinate $N$ has stability $\Theta(2^{-N})$ in the \emph{depth} $N$, because by
\eqref{eq:weight} a $D$-perturbation of size $\delta_0$ moves that coordinate by up to
$2^{N}\delta_0$.  No such factor enters here, because (W1)--(W3) constrain $\rho$ between
two orbits rather than a single deep coordinate of one image.  The finite-horizon layer
is therefore stable precisely where the infinite-depth covering layer is not.
\end{remark}

\begin{definition}[Finite-horizon scrambled packet]\label{def:packet}
Let \(M\ge2\), let \(W_\bullet=(W_n)\), let \(r\ge0\),
\(\varepsilon\in(0,1)\), \(0<\eta<1\), and \(N\ge1\).  We say
\(g\in\F(\Q)\) has the
\emph{\((M,r,\varepsilon,\eta,N)\)-finite-horizon scrambled packet} if there
are points \(x^1,\ldots,x^M\in\Q\), a set
\(P\subseteq\{0,\ldots,N-1\}\), and sets
\(S_{uv}\subseteq\{0,\ldots,N-1\}\), \(1\le u<v\le M\), such that, with
\(G^j_r=g_{r+j}\circ\cdots\circ g_{r+1}\),
\begin{enumerate}
\item[(P1)] \(\#P\ge(1-\eta)N\);
\item[(P2)] \(\#S_{uv}\ge(\eta/8)N\) for every \(1\le u<v\le M\);
\item[(P3)] \(\rho(G^j_r(x^u),G^j_r(x^v))<\varepsilon\) for every
  \(j\in P\) and every \(u<v\);
\item[(P4)] \(\rho(G^j_r(x^u),G^j_r(x^v))>t_0\) for every
  \(j\in S_{uv}\);
\item[(P5)] \(G^j_r(x^u)\in W_{r+j+1}\) for every \(u\le M\) and every
  \(0\le j<N\).
\end{enumerate}
Write
\[
  \mathcal P_{M,r,\varepsilon,\eta,N}(W_\bullet)\subseteq\F(\Q)
\]
for the set of systems with such a packet.
\end{definition}

\begin{lemma}[Finite-packet stability]\label{lem:packet-stability}
If \(g\) has the packet of Definition~\ref{def:packet} with all inequalities strict
by a common margin \(\kappa>0\), then there is \(\delta>0\) such that every
\(h\in\F(\Q)\) with \(D(h,g)<\delta\) has the same packet, with the same
points \(x^u\) and the same sets \(P,S_{uv}\).  In particular
\(\mathcal P_{M,r,\varepsilon,\eta,N}(W_\bullet)\) is open.
\end{lemma}

\begin{proof}
The proof is the same orbit-tracking argument as
Lemma~\ref{lem:fh-stability}.  Let \(\omega\) be a common modulus of continuity
for the finitely many maps \(g_{r+1},\ldots,g_{r+N-1}\).  For each packet point
\(x^u\), the orbit error under a perturbation \(h\) obeys
\[
  e_j^u\le D(h,g)+\omega(e_{j-1}^u),\qquad e_0^u=0 .
\]
Choose \(\delta\) so that all errors are \(<\kappa\) for \(j<N\).  Then every
proximity inequality loses at most \(2\kappa\), every separation inequality
loses at most \(2\kappa\), and every window visit remains inside its open
window.  The counting conditions are unchanged.
\end{proof}

\begin{lemma}[Binary packet codes]\label{lem:packet-code}
For every \(M\ge2\) there is \(L_M\) such that for every \(L\ge L_M\) one can
choose words \(c^1,\ldots,c^M\in\{+,-\}^L\) with pairwise Hamming distance at
least \(L/4\).
\end{lemma}

\begin{proof}
Choose \(M\) independent uniform words of length \(L\).  For a fixed pair, the
Hamming distance has distribution \(\operatorname{Bin}(L,1/2)\).  By the
Chernoff bound,
\[
  \mathbb P\{d_H(c^u,c^v)<L/4\}\le e^{-L/16}.
\]
Thus the probability that some pair violates the \(L/4\) bound is at most
\(\binom M2 e^{-L/16}\), which is \(<1\) once \(L\) is large enough.  Hence a
choice with the desired distances exists.
\end{proof}

\section{Density and the residual theorem}\label{sec:density}

We now show each witness set is dense and assemble Theorem~\ref{thm:A}.  The density
input is Theorem~\ref{thm:dense}; the new ingredient is a block schedule that places
enough proximity blocks and separation blocks before a common horizon to realize the
fraction constraints \(\#P\ge(1-\eta)N\), \(\#S\ge\theta N\), with \(0<\theta<\eta<1\).

\paragraph{Block schedule.}
On the universal schedule of Theorem~\ref{thm:dense} the time axis $\{1,2,\dots\}$ is
partitioned into consecutive blocks $\mathbf B_1,\mathbf B_2,\dots$, where block
$\mathbf B_k$ occupies $L_k$ consecutive times and carries a single label.  We use two
kinds of labelled bricks supplied by the construction of \cite[\S12]{Shmalo2026}:

\begin{itemize}
\item a \emph{proximity brick} $C$ of diameter $<\gamma$, with the property that any two
orbits forced through $C$ at a given time are afterwards $\rho$-distance $<\gamma$ apart
at that time (orbits entering a common small box are close);
\item a \emph{separation brick} $D$ consisting of two sub-bricks \(D^{+},D^{-}\) separated
in the first \(q\) low coordinates, where \(q=q(t_0)\) is any integer with
\(1-2^{-q}>t_0\): the construction of \cite[\S12]{Shmalo2026} supplies, for each such
\(q\), sub-bricks placed in opposite corners of the first \(q\) coordinate intervals, so
that \(|u_i-v_i|=1\) for every coordinate \(i\le q\) whenever \(u\in D^{+}\),
\(v\in D^{-}\).  Then \(\rho(u,v)\ge\sum_{i\le q}2^{-i}|u_i-v_i|=\sum_{i\le q}2^{-i}=
1-2^{-q}>t_0\), so routing one orbit into \(D^{+}\) and the other into \(D^{-}\) forces
\(\rho>t_0\) at that time.
\end{itemize}

Because every label occurs on infinitely many blocks, for any prescribed long horizon we
may select, by the finite-itinerary realizability of Theorem~\ref{thm:dense}, an
admissible itinerary that spends a chosen number of times in proximity bricks and a
chosen number in separation sub-bricks, all inside the prescribed windows.  The only
arithmetic we need is that the block lengths can be made a vanishing fraction of the
horizon, so the fractions \(1-\eta\) and \(\theta\) are attainable with room for the
slack \(\eta-\theta\).

\begin{lemma}[Schedule arithmetic]\label{lem:schedule}
Let \(0<\theta<\eta<1\).  There is a sequence of horizons \(N\to\infty\) and, for each
such \(N\), disjoint sets \(P,S\subseteq\{0,\dots,N-1\}\) with
\(\#P\ge(1-\eta)N\) and \(\#S\ge\theta N\) such that \(P\) is a union of full proximity
blocks and \(S\) is a union of full separation blocks of the schedule (each block placed
at times \(<N\)), provided the block lengths satisfy \(L_k/(L_1+\cdots+L_k)\to0\).
\end{lemma}

\begin{proof}
Choose the schedule so that proximity and separation blocks alternate in long runs.
Write \(T_k=L_1+\cdots+L_k\).  Allocate to \(S\) the separation blocks and to \(P\) the
proximity blocks among \(\mathbf B_1,\dots,\mathbf B_k\), with the run lengths chosen so
that the total separation length \(\sigma_k\) satisfies
\(\theta T_k\le \sigma_k\le\theta T_k+L_k\).  Then
\[
\#S\ge\theta T_k,\qquad \#P\ge(1-\theta)T_k-L_k.
\]
Setting \(N=T_k\), the separation estimate gives \(\#S\ge\theta N\).  Since
\(\theta<\eta\) and \(L_k/T_k\to0\), for all sufficiently large \(k\) we also have
\[
(1-\theta)N-L_k\ge(1-\eta)N.
\]
Thus \(\#P\ge(1-\eta)N\) along \(N=T_k\to\infty\).  Disjointness of \(P,S\) is automatic
since each is a union of distinct whole blocks.
\end{proof}

\begin{lemma}[Density of the witness set]\label{lem:density}
For every \(W_\bullet\), every \(r\ge0\), every \(\varepsilon\in(0,1)\), and every
\(0<\theta<\eta<1\), there are arbitrarily large horizons \(N\) for which
\(\W_{r,\varepsilon,t_0,\eta,\theta,N}(W_\bullet)\) is dense in \((\F(\Q),D)\).
\end{lemma}

\begin{proof}
Fix $f\in\F(\Q)$ and $\eta_0>0$; we find $g$ with $D(f,g)<\eta_0$ in the witness set.
Apply Theorem~\ref{thm:dense} to the window sequence $W_\bullet$ to obtain
$g\in\F(\Q)$, $D(f,g)<\eta_0$, carrying an exact $W_\bullet$-windowed horseshoe on the
universal schedule, with proximity bricks $C$ of diameter $\gamma<\varepsilon$ and
separation bricks $D=D^{+}\cup D^{-}$ separated in the first $q=q(t_0)$ coordinates
(Remark~\ref{rem:t0}).  Fix a tail \(r\ge0\).  By Lemma~\ref{lem:schedule} choose a
horizon \(N\) and disjoint \(P,S\subseteq\{0,\dots,N-1\}\), with
\(\#P\ge(1-\eta)N\) and \(\#S\ge\theta N\), such that \(P\) is a union of proximity
blocks and \(S\) is a union of separation blocks at times \(<N\).

Build the two itineraries from time $r+1$.  Let both orbits visit the \emph{same}
proximity brick at every time $j\in P$, and let one orbit visit $D^{+}$ and the other
$D^{-}$ at every time $j\in S$; at the remaining times $j<N$ (within blocks that are
neither wholly in $P$ nor in $S$, of total length $o(N)$) let both orbits follow any
common admissible brick, which we also count into the windows but into neither $P$ nor
$S$.  This is an admissible finite itinerary of bricks, all lying in their prescribed
windows $W_{r+j+1}$.  By the finite-itinerary realizability of Theorem~\ref{thm:dense}
there exist $x,y\in\Q$ whose $\mathrm{tail}_r$-orbits follow these itineraries exactly
for $0\le j<N$; in particular $G^j_r(x),G^j_r(y)\in W_{r+j+1}$ for all $j<N$, giving
(W3).

For $j\in P$ both orbit points lie in the same proximity brick of diameter
$\gamma<\varepsilon$, so $\rho(G^j_r(x),G^j_r(y))\le\gamma<\varepsilon$, giving (W1).
For $j\in S$ the points lie in $D^{+}$ and $D^{-}$ respectively.  These sub-bricks are
built in \cite[\S12]{Shmalo2026} separated in each of the first $q$ coordinates, so
$|\pi_iG^j_r(x)-\pi_iG^j_r(y)|=1$ for every coordinate index $i\le q$, and therefore
$\rho(G^j_r(x),G^j_r(y))\ge\sum_{i\le q}2^{-i}=1-2^{-q}>t_0$, giving (W2).  Thus \(g\) has
the witness \(\W_{r,\varepsilon,t_0,\eta,\theta,N}(W_\bullet)\) and
$D(f,g)<\eta_0$.  Since $f,\eta_0$ were arbitrary, the witness set is dense.
\end{proof}

\begin{remark}[Calibrating $t_0$ to the separation bricks]\label{rem:t0}
The separation scale $t_0$ in Definition~\ref{def:witness} is fixed in advance.  A single
first-coordinate gap contributes at most $2^{-1}\cdot1=\tfrac12$ to $\rho$, so separating
in coordinate $1$ alone caps the achievable separation at $\tfrac12$ and only handles
$t_0<\tfrac12$.  To reach any $t_0\in(0,1)$ we separate in the first $q$ low coordinates:
fix an integer $q=q(t_0)$ with $1-2^{-q}>t_0$, and use the separation bricks of
\cite[\S12]{Shmalo2026} whose sub-bricks $D^{+},D^{-}$ sit in opposite corners of the
first $q$ coordinate intervals (the $\dist_q$ construction there supplies exactly these),
so that $\rho\ge\sum_{i\le q}2^{-i}=1-2^{-q}>t_0$ at separation times.  Fully separating
the first $q$ coordinates reaches $1-2^{-q}$, which exceeds any $t_0<1$ once $q$ is large
enough.  The strict-witness hypothesis of Lemma~\ref{lem:fh-stability} is met because
$1-2^{-q}>t_0$ by a positive margin and the proximity diameter is $<\varepsilon$ by a
positive margin, with window points in the interiors of the (open) windows; one fixes
$\kappa>0$ smaller than all these margins.
\end{remark}

\begin{lemma}[Density of finite scrambled packets]\label{lem:packet-density}
For every \(M\ge2\), every \(W_\bullet\), every \(r\ge0\), every
\(\varepsilon\in(0,1)\), and every \(0<\eta<1\), there are arbitrarily large
horizons \(N\) for which
\(\mathcal P_{M,r,\varepsilon,\eta,N}(W_\bullet)\) is dense in \((\F(\Q),D)\).
\end{lemma}

\begin{proof}
Fix \(f\in\F(\Q)\) and \(\eta_0>0\).  Apply Theorem~\ref{thm:dense} to get
\(g\) with \(D(f,g)<\eta_0\) carrying an exact windowed horseshoe.  Choose the
proximity bricks with diameter \(\gamma<\varepsilon\), and choose separation
bricks \(D^+,D^-\) calibrated to \(t_0\) as in Remark~\ref{rem:t0}.

Use Lemma~\ref{lem:schedule} with the pair \((\theta,\eta)\) replaced by
\((\eta/2,\eta)\).  Along arbitrarily large horizons \(N\) this gives disjoint
sets \(P,S\subseteq\{0,\ldots,N-1\}\) such that \(P\) is a union of proximity
blocks, \(S\) is a union of separation blocks,
\[
  \#P\ge(1-\eta)N,\qquad \#S\ge(\eta/2)N .
\]
Taking \(N\) still larger if necessary, assume \(\#S\ge L_M\), where
\(L_M\) is from Lemma~\ref{lem:packet-code}.  Choose words
\[
  c^1,\ldots,c^M\in\{+,-\}^{S}
\]
with pairwise Hamming distance at least \(\#S/4\).

Build \(M\) finite itineraries from time \(r+1\).  At every time \(j\in P\), all
\(M\) itineraries visit the same proximity brick.  At every separation time
\(j\in S\), the \(u\)-th itinerary visits \(D^+\) or \(D^-\) according to the
symbol \(c^u_j\).  At all remaining times below \(N\), all itineraries follow
any common admissible brick.  Finite-itinerary realizability gives points
\(x^1,\ldots,x^M\) whose tail-\(r\) orbits follow these itineraries.

All orbit points lie in the prescribed windows.  On \(P\), all points lie in
the same proximity brick, so all pairwise distances are \(<\varepsilon\).  For
each pair \(u<v\), let
\[
  S_{uv}:=\{j\in S:\ c^u_j\ne c^v_j\}.
\]
Then \(\#S_{uv}\ge \#S/4\ge(\eta/8)N\), and at every \(j\in S_{uv}\) the two
orbits lie in opposite separation sub-bricks, hence are \(t_0\)-separated.
Thus \(g\in\mathcal P_{M,r,\varepsilon,\eta,N}(W_\bullet)\), proving density.
\end{proof}

\begin{theorem}[Residual finite-horizon windowed DC1 witnesses]\label{thm:residual}
Let $\{W^a_\bullet\}_{a\ge1}$ be a countable family of sequences of dense open subsets of
$\Q$ and let \(t_0\in(0,1)\).  Then
\[
\Gcal\;=\;\bigcap_{a\ge1}\ \bigcap_{r\ge0}\ \bigcap_{\substack{\varepsilon\in\mathbb Q\\0<\varepsilon<1}}\ \bigcap_{\substack{\eta,\theta\in\mathbb Q\\0<\theta<\eta<1}}\ \bigcap_{N_0\ge1}\ \bigcup_{N\ge N_0}\W_{r,\varepsilon,t_0,\eta,\theta,N}(W^a_\bullet)
\]
is comeager in $(\F(\Q),D)$, and every $g\in\Gcal$ satisfies the conclusion of
Theorem~\ref{thm:A}.
\end{theorem}

\begin{proof}
For fixed \(a,r,\varepsilon,\eta,\theta,N_0\) the set
\[
U^{a,r,\varepsilon,\eta,\theta}_{N_0}
=\bigcup_{N\ge N_0}\W_{r,\varepsilon,t_0,\eta,\theta,N}(W^a_\bullet)
\]
is open by Lemma~\ref{lem:fh-stability} (a union of open sets) and dense: by
Lemma~\ref{lem:density} there are arbitrarily large horizons $N$ for which
\(\W_{r,\varepsilon,t_0,\eta,\theta,N}(W^a_\bullet)\) is dense, so in particular some
\(N\ge N_0\) gives a dense subset of \(U^{a,r,\varepsilon,\eta,\theta}_{N_0}\), whence
the union is dense.  The index set of all \(a,r,N_0\in\mathbb N\), rational
\(\varepsilon\in(0,1)\), and rational \(0<\theta<\eta<1\) is countable, so \(\Gcal\) is
a countable intersection of open dense sets, hence comeager
by Lemma~\ref{lem:baire}.  If $g\in\Gcal$ then for every $a,r$, every rational
\(\varepsilon\in(0,1)\), and all rational \(0<\theta<\eta<1\), for every \(N_0\) there
is \(N\ge N_0\) with
\(g\in\W_{r,\varepsilon,t_0,\eta,\theta,N}(W^a_\bullet)\), i.e.\ there are
$x,y,P,S$ as in Definition~\ref{def:witness} for that $N$; this is exactly the statement
of Theorem~\ref{thm:A} (for irrational $\varepsilon$ pass to a smaller rational, which
only strengthens (W1); the tolerances \(\eta,\theta\) are required rational, as stated).
\end{proof}

\begin{theorem}[Residual finite-horizon scrambled packets]\label{thm:packet}
Let $\{W^a_\bullet\}_{a\ge1}$ be a countable family of sequences of dense open subsets
of $\Q$ and let \(t_0\in(0,1)\).  Then
\[
\Gcal_{\rm pkt}
\;=\;
\bigcap_{a\ge1}\ \bigcap_{r\ge0}\ \bigcap_{M\ge2}\
\bigcap_{\substack{\varepsilon\in\mathbb Q\\0<\varepsilon<1}}\
\bigcap_{\substack{\eta\in\mathbb Q\\0<\eta<1}}\
\bigcap_{N_0\ge1}\ \bigcup_{N\ge N_0}
\mathcal P_{M,r,\varepsilon,\eta,N}(W^a_\bullet)
\]
is comeager in \((\F(\Q),D)\).  Every \(g\in\Gcal_{\rm pkt}\) satisfies the
packet conclusion stated in Theorem~\ref{thm:packet-intro}.
\end{theorem}

\begin{proof}
For fixed \(a,r,M,\varepsilon,\eta,N_0\), the union over \(N\ge N_0\) is open by
Lemma~\ref{lem:packet-stability} and dense by Lemma~\ref{lem:packet-density}.
The displayed intersection is countable, hence comeager.  Membership gives
arbitrarily large horizons with the packet of Definition~\ref{def:packet}.  Passing
from rational \(\varepsilon,\eta\) to the stated rational parameters is immediate.
\end{proof}

\begin{remark}[What the witnesses approximate]\label{rem:approx}
Fix $g\in\Gcal$, a tail $r$ and a scale $\varepsilon$.  Taking
\(\eta=1/m\) and \(\theta=1/(2m)\) for \(m=2,3,\dots\) yields, for each \(m\),
arbitrarily long horizons carrying a pair proximal on a \((1-1/m)\)-fraction of times
and \(t_0\)-separated on at least a \(1/(2m)\)-fraction.  This
is the finite-horizon, fixed-tolerance silhouette of a DC1 scrambled pair: the upper
proximity statistic tends to $1$ and the separation statistic stays positive at the
fixed scale $t_0$, with all visits inside the prescribed windows up to the horizon.  We
do not assert a single pair valid for all horizons at once; producing such a pair is the
infinite-horizon problem (Problem~\ref{prob:companion}), which Section~\ref{sec:obstruction}
shows no finite-horizon certificate can reach.
\end{remark}

\begin{remark}[Finite packets are still finite-horizon]\label{rem:packet-honest}
Theorem~\ref{thm:packet} strengthens the positive side from one pair to
arbitrarily large finite packets.  It is still a finite-horizon statement: the
packet points and the code across the separation times may depend on \(N\), and
the theorem does not produce one uncountable scrambled set, or even one finite
packet, working along all horizons simultaneously.  The binary coding only says
that the dense horseshoe's two separation sub-bricks can be reused over many
separation times to make every pair in a finite list disagree often.
\end{remark}

\section{The obstruction: no finite-horizon certificate suffices}\label{sec:obstruction}

We now prove Theorem~\ref{thm:B}.  We first recall metric mean dimension, since the
squeeze argument controls it directly and infinite topological entropy is one of its
consequences in the relevant direction.

For a non-autonomous system $h$ on $\Q$, a tail $\mathrm{tail}_sh$, scale
$\varepsilon>0$ and horizon $N$, let $\mathrm{sep}_{s,N}(h,\varepsilon)$ be the maximal
cardinality of a set $E\subseteq\Q$ that is $(s,N,\varepsilon)$-separated in the Bowen
metric $\rho^h_{s,N}(x,y)=\max_{0\le j<N}\rho(H^j_s(x),H^j_s(y))$, and put
\[
\mdim(\mathrm{tail}_sh)
=\limsup_{\varepsilon\downarrow0}\ \frac{1}{\log(1/\varepsilon)}\
\limsup_{N\to\infty}\frac1N\log\mathrm{sep}_{s,N}(h,\varepsilon).
\]
The inner double limit is the topological entropy $h_{\mathrm{top}}(\mathrm{tail}_sh)$
at scale $\varepsilon$; if $h_{\mathrm{top}}(\mathrm{tail}_sh)=\infty$ then
$\frac1N\log\mathrm{sep}_{s,N}\to\infty$ as $\varepsilon\downarrow0$ along the inner
limit, and one checks $\mdim$ may still be finite or infinite.  We only use the
following one-directional bound.

\begin{lemma}[Squeeze cap]\label{lem:cap}
Let $L\ge1$ and let $h\in\F(\Q)$ satisfy $\pi_j h_n(x)=\tfrac12$ for every $n$, every
$x\in\Q$, and every coordinate $j>L$ --- that is, every image of every $h_n$ lies in the
$L$-dimensional slice $Z_L=\{z\in\Q:\ z_j=\tfrac12\ \forall j>L\}$.  Then for every tail
$s\ge0$ and every fixed scale \(0<\varepsilon\le1\),
\begin{equation}\label{eq:entscale}
h_{\mathrm{top}}^{(\varepsilon)}(\mathrm{tail}_sh)
:=\limsup_{N\to\infty}\frac1N\log\mathrm{sep}_{s,N}(h,\varepsilon)
\ \le\ L\log(4/\varepsilon),
\end{equation}
so the entropy at every fixed scale is finite and $\mdim(\mathrm{tail}_sh)\le L<\infty$.
\end{lemma}

\begin{proof}
Fix $s\ge0$.  Since every $h_n$ maps into $Z_L$, for every $i\ge1$ and every $x$ the
point $H^i_s(x)$ lies in $Z_L$; only the base point $H^0_s(x)=x$ can lie outside $Z_L$.
The slice $Z_L$ is isometric, in the metric $\rho$ restricted to it, to $[0,1]^L$ with the
weighted metric $d_L(u,v)=\sum_{j\le L}2^{-j}|u_j-v_j|$, because coordinates $j>L$ are
constant on $Z_L$.  A grid of mesh $\varepsilon/2$ in each of the first $L$ coordinates
has $(\lceil 2/\varepsilon\rceil)^{L}$ cells and resolves $d_L$ to within
$\sum_{j\le L}2^{-j}(\varepsilon/2)<\varepsilon/2<\varepsilon$; hence the $\rho$-covering
number of $Z_L$ at radius $\varepsilon$ satisfies
\[
\mathrm{cov}_\rho(Z_L,\varepsilon)\ \le\ (\lceil 2/\varepsilon\rceil)^{L}\ \le\ (4/\varepsilon)^{L}
\qquad(\varepsilon\le1).
\]

Now bound the Bowen separation.  An $(s,N,\varepsilon)$-separated set $E$ is, in
particular, mapped injectively by $x\mapsto(H^0_s(x),H^1_s(x),\dots,H^{N-1}_s(x))$ to a
set that is $\varepsilon$-separated in the max-$\rho$ metric on the product
$\Q\times Z_L\times\cdots\times Z_L$ ($N-1$ copies of $Z_L$).  An $\varepsilon$-separated
set in a product (max metric) has cardinality at most the product of the factors'
$\varepsilon$-covering numbers, so
\[
\mathrm{sep}_{s,N}(h,\varepsilon)\ \le\ \mathrm{cov}_\rho(\Q,\varepsilon)\cdot
\mathrm{cov}_\rho(Z_L,\varepsilon)^{\,N-1}
\ \le\ \mathrm{cov}_\rho(\Q,\varepsilon)\cdot(4/\varepsilon)^{L(N-1)} .
\]
The factor $\mathrm{cov}_\rho(\Q,\varepsilon)$ is a finite constant depending only on
$\varepsilon$ (the Hilbert cube is compact), and crucially it is independent of $N$.
Therefore
\[
\frac1N\log\mathrm{sep}_{s,N}(h,\varepsilon)
\ \le\ \frac{\log\mathrm{cov}_\rho(\Q,\varepsilon)}{N}
+\frac{N-1}{N}\,L\log(4/\varepsilon)\ \xrightarrow[N\to\infty]{}\ L\log(4/\varepsilon).
\]
Thus \eqref{eq:entscale} holds, the right-hand side is finite for each fixed
\(0<\varepsilon\le1\),
and dividing by $\log(1/\varepsilon)$ and letting $\varepsilon\downarrow0$,
\[
\mdim(\mathrm{tail}_sh)
=\limsup_{\varepsilon\downarrow0}\frac{h^{(\varepsilon)}_{\mathrm{top}}(\mathrm{tail}_sh)}{\log(1/\varepsilon)}
\ \le\ \limsup_{\varepsilon\downarrow0}\frac{L\log(4/\varepsilon)}{\log(1/\varepsilon)}=L. \qedhere
\]
\end{proof}

\begin{remark}[Why deep coordinates are collapsed, not merely confined]\label{rem:collapse}
It is essential that the squeeze sends deep coordinates to a single value rather than
into a proper subinterval.  A perturbation confining coordinates $j>L$ to a band
$[\frac18,\frac78]$ would leave those coordinates free to vary, and the band
$\prod_{j>L}[\frac18,\frac78]$ is still an infinite-dimensional Hilbert-cube tail: a map
into it can have arbitrarily large metric mean dimension, so confinement to a band does
\emph{not} cap $\mdim$.  Collapsing coordinates $j>L$ to the constant $\frac12$ lands
every image in the genuinely $L$-dimensional slice $Z_L$, of finite box dimension $L$,
which is what forces \eqref{eq:entscale}.  The $D$-cost is the same to leading order:
$\sum_{j>L}2^{-j}|y_j-\frac12|\le\frac12 2^{-L}$.
\end{remark}

\begin{remark}[Which complexity notion the cap contradicts]\label{rem:notion}
We phrase the obstruction through metric mean dimension, which is unambiguous across
non-autonomous entropy conventions: \eqref{eq:entscale} gives
$\mdim(\mathrm{tail}_sh)\le L<\infty$ for every $s$.  The ``infinite topological
entropy'' tails produced by \cite{Shmalo2026} satisfy the stronger scale law
$h^{(\varepsilon)}_{\mathrm{top}}(\mathrm{tail}_sg)\ge q\log(1/\varepsilon)-C_q$ for all
$q$, equivalently $\mdim=\infty$; that is the property the squeeze defeats, since
\eqref{eq:entscale} caps the same scale-rate by \(L\log(4/\varepsilon)\), which cannot
exceed $q\log(1/\varepsilon)-C_q$ for all small $\varepsilon$ once $q>L$.  Thus
Theorem~\ref{thm:obstruction} is proved for the metric-mean-dimension form, and the
entropy form follows from the same inequality.
\end{remark}

\begin{theorem}[Finite-horizon certificate obstruction]\label{thm:obstruction}
Let $V\subseteq\F(\Q)$ be finite-horizon certified: there are $r\ge0$ and $N\ge1$ such
that whether $g\in V$ depends only on $(g_{r+1},\dots,g_{r+N})$.  Suppose
\[
V\subseteq\bigl\{\,g\in\F(\Q):\ \mdim(\mathrm{tail}_sg)=\infty\ \text{for some }s\ge0\,\bigr\}.
\]
Then $V=\varnothing$.  The same conclusion holds if $V$ is contained in the class of
systems some tail of which has infinite topological entropy in the scale-law sense of
\cite{Shmalo2026} (i.e.\ $h^{(\varepsilon)}_{\mathrm{top}}\ge q\log(1/\varepsilon)-C_q$
for all $q$).
\end{theorem}

\begin{proof}
Suppose $V\ne\varnothing$ and fix $g\in V$.  Since $V$ is open and finite-horizon
certified by the window $(g_{r+1},\dots,g_{r+N})$, there is $\delta>0$ such that every
$h\in\F(\Q)$ with $h_{r+i}=g_{r+i}$ for $1\le i\le N$ and $D(h,g)<\delta$ lies in $V$;
indeed, membership ignores all coordinates outside the window, so any $h$ agreeing with
$g$ on the window is in $V$ regardless of its other coordinates, provided it is within
$D$-distance $\delta$ of $g$ where $\delta$ is the openness radius (which only constrains
the windowed coordinates, on which $h$ equals $g$, so any such $h$ qualifies).

Choose $L\ge1$ with $2^{-L}<\delta$.  Define the coordinate squeeze $\Phi:\Q\to\Q$,
collapsing every coordinate past depth $L$ to the constant $\tfrac12$:
\[
\Phi(y)_j=
\begin{cases}
y_j, & j\le L,\\[2pt]
\tfrac12, & j>L.
\end{cases}
\]
$\Phi$ is continuous, $\Phi(\Q)=Z_L:=\{z\in\Q:\ z_j=\tfrac12\ \forall j>L\}$, and
\[
\rho(\Phi(y),y)=\sum_{j>L}2^{-j}\Bigl|y_j-\tfrac12\Bigr|
\le\sum_{j>L}2^{-j}\cdot\tfrac12=\tfrac12\,2^{-L}<\delta .
\]
Define $h\in\F(\Q)$ by leaving the certifying window untouched and squeezing every other
map:
\[
h_n=\begin{cases}
g_n, & n\in\{r+1,\dots,r+N\},\\
\Phi\circ g_n, & \text{otherwise.}
\end{cases}
\]
Then $d_\infty(h_n,g_n)=0$ on the window and
$d_\infty(h_n,g_n)=\sup_x\rho(\Phi(g_n(x)),g_n(x))<\delta$ off the window, so
$D(h,g)<\delta$, and $h$ agrees with $g$ on the certifying window; hence $h\in V$.

Now examine an arbitrary tail of $h$.  Fix $s\ge0$ and consider the orbit
$H^i_s=h_{s+i}\circ\cdots\circ h_{s+1}$.  For $i\ge1$ the outermost map is $h_{s+i}$,
which is either a squeezed map $\Phi\circ g_{s+i}$ --- in which case $H^i_s(x)\in Z_L$
for every $x$ --- or, if $s+i\in\{r+1,\dots,r+N\}$, the untouched map $g_{s+i}$.  The
window is a set of at most $N$ times, so the second case occurs for at most $N$ values of
$i\ge1$; call this exceptional set $\mathcal E\subseteq\{1,\dots\}$, $\#\mathcal E\le N$.
For all $i\ge1$ with $i\notin\mathcal E$ we have $H^i_s(x)\in Z_L$ for every $x$.

Bound the Bowen separation exactly as in Lemma~\ref{lem:cap}, isolating the finitely
many exceptional time-slices.  An $(s,N',\varepsilon)$-separated set injects, via
$x\mapsto(H^i_s(x))_{0\le i<N'}$, into a set that is $\varepsilon$-separated in the
max-$\rho$ metric on $\prod_{0\le i<N'}Y_i$, where $Y_i=Z_L$ for $i\ge1$, $i\notin\mathcal E$,
and $Y_i=\Q$ for $i=0$ or $i\in\mathcal E$.  Hence
\[
\mathrm{sep}_{s,N'}(h,\varepsilon)\ \le\ \mathrm{cov}_\rho(\Q,\varepsilon)^{\,1+\#\mathcal E}\cdot
\mathrm{cov}_\rho(Z_L,\varepsilon)^{\,N'-1-\#\mathcal E}
\ \le\ \mathrm{cov}_\rho(\Q,\varepsilon)^{\,1+N}\cdot(4/\varepsilon)^{L N'} .
\]
The first factor is independent of $N'$ (the window length $N$ is fixed), so dividing by
$N'$ and letting $N'\to\infty$ gives, as in \eqref{eq:entscale},
\[
h^{(\varepsilon)}_{\mathrm{top}}(\mathrm{tail}_sh)\ \le\ L\log(4/\varepsilon),
\qquad
\mdim(\mathrm{tail}_sh)\ \le\ L<\infty,
\]
for every $s\ge0$.  This contradicts
$h\in V\subseteq\{\,\mdim(\mathrm{tail}_sg)=\infty\ \text{for some }s\,\}$.  Hence
$V=\varnothing$.

For the scale-law entropy form, an $h\in V$ would need some tail with
$h^{(\varepsilon)}_{\mathrm{top}}\ge q\log(1/\varepsilon)-C_q$ for all $q$; but the bound
$h^{(\varepsilon)}_{\mathrm{top}}(\mathrm{tail}_sh)\le L\log(4/\varepsilon)$ fails this
for $q>L$ at all small $\varepsilon$.  So $V=\varnothing$ in that case as well.
\end{proof}

\begin{remark}[A horizon obstruction, complementary to the depth obstruction]\label{rem:horizon}
The squeeze in Theorem~\ref{thm:obstruction} acts on all maps \emph{outside a fixed
window of times}, whereas the empty-interior squeeze of \cite[Theorem~6.2]{companion}
acts on all coordinates \emph{past a fixed depth} of a single map.  The first defeats
finite-\emph{horizon} certificates of infinite-horizon complexity; the second defeats
open sets requiring infinite-\emph{depth} covering.  Both rest on the same elementary
fact --- a coordinate (or time) tail of total $D$-weight $<\delta$ can be overwritten
without leaving a $\delta$-ball --- but they constrain orthogonal method classes.  Read
together: no open set can certify infinite depth (depth obstruction), and no
finite-horizon-certified open set can certify an infinite-\(\mdim\) tail or a scale-law
entropy tail (horizon obstruction).
\end{remark}

\begin{corollary}[Banach--Mazur form]\label{cor:bm}
Let $V$ be finite-horizon certified and contained in one of the complexity classes covered
by Theorem~\ref{thm:obstruction}.  Then in the Banach--Mazur game on
$(\F(\Q),D)$, Player~II (the player choosing the even moves) can avoid $V$ entirely:
after any move of Player~I, Player~II plays a ball produced by the squeeze of
Theorem~\ref{thm:obstruction}, all of whose points lie outside $V$.  Equivalently,
$V$ is not residual; being also open, $V$ is empty.  Consequently no countable
union of finite-horizon-certified open sets, each entailing an infinite-\(\mdim\) tail or
one of the scale-law entropy tails of Theorem~\ref{thm:obstruction}, can be dense, and the
Baire strategy through such a family cannot prove Problem~\ref{prob:companion}.
\end{corollary}

\begin{proof}
By Theorem~\ref{thm:obstruction}, $V=\varnothing$, so every $D$-ball is disjoint from
$V$; in particular Player~II has the trivial avoidance strategy, and the open empty set
is vacuously non-residual.  The final clause follows because a countable union of empty
sets is empty, hence not dense, so its members cannot serve as the dense-open building
blocks of a Baire argument.
\end{proof}

\subsection*{The squeeze also defeats semifactors and confines scrambling}

The cap of Lemma~\ref{lem:cap} is an entropy statement, and Theorem~\ref{thm:obstruction}
phrases the obstruction through $\mdim$ and the scale-law entropy.  The same squeezed
system $h$ fails a further growth-type conclusion of the dense theorem outright, with no
horizon hypothesis: the all-tail symbolic semifactor onto an unbounded-alphabet shift.
This sharpens what the method defeats and, by contrast, isolates the one conclusion it
leaves standing.

We use the following operational notion, which is the only feature of a symbolic
semifactor onto a shift we need.  Say a tail $\mathrm{tail}_sh$ admits a \emph{compact
symbolic semifactor with branch separation $\varepsilon_*>0$ and block alphabets
$(M_n)_{n\ge1}$} if for every $n$ the system carries an $(s,n,\varepsilon_*)$-separated
set of cardinality at least $\prod_{i\le n}M_i$.  This holds whenever $\mathrm{tail}_sh$
factors onto a one-sided shift whose length-$n$ admissible words number
$\prod_{i\le n}M_i$ and whose distinct length-$n$ cylinders pull back to orbit segments
lying in branch boxes pairwise more than $\varepsilon_*$ apart in the depth-$n$ Bowen
metric $\rho^h_{s,n}$: the factor map sends each length-$n$ Bowen ball into a single
cylinder, so distinct admissible words are realized by $\varepsilon_*$-separated segments.
The full variable-alphabet shift of Theorem~\ref{thm:dense}, with $M_n\to\infty$, is the
case of interest.

\begin{lemma}[Semifactor obstruction]\label{lem:semifactor}
Let $L\ge1$ and let $h\in\F(\Q)$ have every image of every $h_n$ in the slice
$Z_L=\{z\in\Q:\ z_j=\tfrac12\ \forall j>L\}$.  Then no tail $\mathrm{tail}_sh$ admits a
compact symbolic semifactor with any fixed branch separation $\varepsilon_*>0$ and block
alphabets $M_n\to\infty$; in particular no tail factors onto the full variable-alphabet
shift of Theorem~\ref{thm:dense}.
\end{lemma}

\begin{proof}
Suppose some tail $\mathrm{tail}_sh$ did, with separation $\varepsilon_*$ and alphabets
$M_n\to\infty$.  By definition the system carries, for each $n$, an
$(s,n,\varepsilon_*)$-separated set of cardinality at least $\prod_{i\le n}M_i$, so
\[
\mathrm{sep}_{s,n}(h,\varepsilon_*)\ \ge\ \prod_{i\le n}M_i,
\qquad
\frac{1}{n}\log\mathrm{sep}_{s,n}(h,\varepsilon_*)\ \ge\
\frac{1}{n}\sum_{i\le n}\log M_i .
\]
The right-hand side is the Cesàro mean of the sequence $(\log M_i)_i$, which tends to
$\infty$ because $M_i\to\infty$; hence
$h^{(\varepsilon_*)}_{\mathrm{top}}(\mathrm{tail}_sh)
=\limsup_{N\to\infty}\frac1N\log\mathrm{sep}_{s,N}(h,\varepsilon_*)=\infty$.
But Lemma~\ref{lem:cap} applies to $h$ verbatim and gives
$h^{(\varepsilon_*)}_{\mathrm{top}}(\mathrm{tail}_sh)\le L\log(4/\varepsilon_*)<\infty$
at the fixed scale $\varepsilon_*$, a contradiction.  Hence no such semifactor exists.
The variable-alphabet shift of Theorem~\ref{thm:dense} has $M_n\to\infty$ and fixed branch
separation, so it is excluded.
\end{proof}

In the obstruction proof of Theorem~\ref{thm:obstruction} the squeezed map agrees with
$g$ on the certifying window and squeezes elsewhere, so only finitely many tail
time-slices escape $Z_L$.  The finite-exception bookkeeping that turned
\eqref{eq:entscale} into the in-proof cap applies here too: adjoining the fixed finite set
$\mathcal E$ of unconstrained Bowen coordinates multiplies $\mathrm{sep}_{s,n}$ by a
constant $\le\mathrm{cov}_\rho(\Q,\varepsilon_*)^{1+\#\mathcal E}$, which changes
$\frac1n\log\mathrm{sep}_{s,n}$ by $o(1)$.  So Lemma~\ref{lem:semifactor} holds for that
$h$ as well, and the squeezed system carries no unbounded-alphabet semifactor on any tail.

\begin{proposition}[Finite-dimensional confinement of scrambling]\label{prop:confine}
Let $L\ge1$ and let $h\in\F(\Q)$ have every image of every $h_n$ in $Z_L$.  Fix a tail
$s\ge0$.  For every $x\in\Q$ the orbit $H^i_s(x)$ lies in $Z_L$ for all $i\ge1$;
consequently every DC1-scrambled set of $\mathrm{tail}_sh$ has its entire post-initial
orbit data inside the $L$-dimensional slice $Z_L$, and the proximal/distal distributions
that define scrambling are computed from orbits in $Z_L$.
\end{proposition}

\begin{proof}
For $i\ge1$ the outermost map of $H^i_s=h_{s+i}\circ\cdots\circ h_{s+1}$ is $h_{s+i}$,
whose image lies in $Z_L$; hence $H^i_s(x)\in Z_L$ for every $x$ and every $i\ge1$.  Only
the base point $H^0_s(x)=x$ can lie outside $Z_L$.  The distributional functions
$F_{xy}$, $F^*_{xy}$ that define a DC1-scrambled pair are limits of empirical
distributions of the distances $\rho(H^j_s(x),H^j_s(y))$ over $0\le j<N$ as
$N\to\infty$; the single term $j=0$ contributes weight $1/N\to0$, so these distributions
are determined by the $Z_L$-valued part of the orbits.  Thus the scrambling lives in
$Z_L$.
\end{proof}

\begin{remark}[The squeeze does \emph{not} kill a bare DC1 set]\label{rem:dc1-honest}
Proposition~\ref{prop:confine} is a confinement statement, not a destruction statement,
and it is worth being exact about the difference.  Collapsing coordinates $j>L$ leaves the
first $L$ coordinates --- in particular coordinate $1$ --- untouched, and the
$t_0$-separation that drives the horseshoe is read in the low coordinates $i\le q\le L$
(Remark~\ref{rem:t0}); the squeeze never touches them.  Uncountable DC1 scrambling and
continuum cardinality are insensitive to ambient dimension: the one-dimensional slice
$Z_1\cong[0,1]$ already carries perfect (hence continuum-sized) DC1-scrambled sets for
suitable interval maps, and confining a scrambled set to $Z_L$ neither reduces its
cardinality nor disturbs its proximal/distal statistics.  So the coordinate squeeze, which
caps entropy and kills the unbounded-alphabet semifactor, does \emph{not} by itself defeat
the existence of an uncountable DC1-scrambled set: that conclusion is dimension-stable and
survives the collapse.  What is proved is only that any such set is finite-dimensionally
confined after one step.
\end{remark}

\begin{corollary}[The conjunction fails on the growth conjuncts only]\label{cor:bundle}
Let $V\subseteq\F(\Q)$ be finite-horizon certified and contained in the class of systems
some tail of which simultaneously (a) carries an uncountable windowed DC1-scrambled set,
(b) has infinite topological entropy in the scale-law sense, and (c) admits a compact
symbolic semifactor onto the full variable-alphabet shift.  Then $V=\varnothing$.  The
emptiness is forced by conjuncts \emph{(b)} and \emph{(c)} alone: the squeeze of
Theorem~\ref{thm:obstruction} produces, in every $D$-ball meeting the certifying window, a
system whose tails violate both (b) (by Lemma~\ref{lem:cap}) and (c) (by
Lemma~\ref{lem:semifactor}), while leaving conjunct (a) attainable in the confined slice
$Z_L$ (Proposition~\ref{prop:confine}, Remark~\ref{rem:dc1-honest}).
\end{corollary}

\begin{proof}
If $V\ne\varnothing$, pick $g\in V$ and form the squeezed $h$ of
Theorem~\ref{thm:obstruction}: $h$ agrees with $g$ on the certifying window,
$D(h,g)<\delta$, so $h\in V$, yet every tail of $h$ has
$h^{(\varepsilon)}_{\mathrm{top}}\le L\log(4/\varepsilon)$ (Lemma~\ref{lem:cap} with the
finite-exception bookkeeping of the theorem) and admits no $M_n\to\infty$ semifactor
(Lemma~\ref{lem:semifactor}, likewise).  Either failure contradicts $h\in V$, since
membership requires (b) and (c) of some tail.  Hence $V=\varnothing$.  That the
obstruction does not pass through (a) is Remark~\ref{rem:dc1-honest}: the same $h$ can
still carry an uncountable scrambled set, now confined to $Z_L$.
\end{proof}

\section{The ledger and what remains open}\label{sec:open}

We restate the open problem of \cite{companion} (which itself isolates
\cite[\S14]{Shmalo2026}) and record precisely what the present results add.

\begin{problem}[{Residuality of the infinite-horizon conclusions; \cite[Problem~8.1]{companion}}]\label{prob:companion}
Let $W_\bullet$ be a sequence of dense open subsets of $\Q$ and let
$\mathcal C(W_\bullet)$ be the class of $g\in\F(\Q)$ such that for every $r\ge0$ the tail
$\mathrm{tail}_rg$ admits a compact symbolic semifactor onto a full variable-alphabet
shift, a continuum-sized DC1-scrambled set visiting $W_{r+j+1}$ at tail time $j$,
infinite topological entropy, and infinite metric mean dimension.  By
\cite[Theorem~3.6]{Shmalo2026}, $\mathcal C(W_\bullet)$ is dense.  Is
$\mathcal C(W_\bullet)$ residual?
\end{problem}

\paragraph{Ledger.}
\begin{itemize}
\item[(i)] \emph{Positive, this paper.}  The finite-horizon windowed DC1 witnesses are
residual: Theorem~\ref{thm:residual} produces a comeager $\Gcal$ on which, at every
tail, every rational scale, every rational \(0<\theta<\eta<1\) and the fixed separation scale
\(t_0\), arbitrarily long witness horizons exist.  This is the finite-horizon shadow of the DC1
half of $\mathcal C(W_\bullet)$, with $\eta\downarrow0$ recovering the proximity
statistic and $t_0$ the separation statistic.  Theorem~\ref{thm:packet} strengthens
this from one pair to arbitrary finite packets: for every \(M\), every tail and every
proximal scale there are arbitrarily long horizons with \(M\) windowed orbit segments
which are mutually proximal on a common \(1-\eta\) fraction and whose every pair is
\(t_0\)-separated on at least an \(\eta/8\) fraction.
\item[(ii)] \emph{Negative, growth-type conclusions.}  No finite-horizon certificate
forces the \emph{growth-type} infinite-horizon conclusions.  Theorem~\ref{thm:obstruction}
shows every finite-horizon-certified open set contained in the infinite-\(\mdim\) class,
or in the scale-law entropy class from \cite{Shmalo2026}, is empty, and the same squeeze
also kills the unbounded-alphabet semifactor (Lemma~\ref{lem:semifactor}); the bundled
Corollary~\ref{cor:bundle} records that the conjunction of entropy, semifactor and DC1
fails on its entropy and semifactor conjuncts alone, while Corollary~\ref{cor:bm} renders
the natural Baire strategy through such certificates impossible.
\item[(ii$'$)] \emph{What the squeeze does \emph{not} defeat.}  The method does not defeat
the bare existence of an uncountable DC1-scrambled set.  By
Proposition~\ref{prop:confine} the squeeze only confines scrambling to a finite-dimensional
slice $Z_L$, and DC1 scrambling with continuum cardinality survives such confinement
(Remark~\ref{rem:dc1-honest}).  So the coordinate-collapse certificate obstruction
separates cleanly: it defeats the entropy/mean-dimension/semifactor (growth-rate)
conclusions, but leaves the DC1 cardinality conclusion open to finite-horizon methods, and
a different obstruction --- or a positive construction --- would be needed for that
conjunct.
\item[(iii)] \emph{Still open.}  Whether $\mathcal C(W_\bullet)$ is residual remains
open.  Theorem~\ref{thm:obstruction} does \emph{not} prove $\mathcal C(W_\bullet)$
non-comeager: a comeager set need not be a countable intersection of finite-horizon
certified open sets, and an infinite-horizon property may be comeager while admitting no
finite-horizon certificate at all (it would be witnessed by an unbounded, system-dependent
window).  We make no claim either way.
\end{itemize}

\paragraph{Where a positive proof must live.}
Combining (ii) with the depth obstruction of \cite[Theorem~6.2]{companion}: a residual
proof of Problem~\ref{prob:companion} cannot use open building blocks that certify
either infinite depth or an infinite-\(\mdim\)/scale-law entropy tail from a bounded
window of times.  The
finite-horizon witnesses of Theorem~\ref{thm:residual} are exactly the strongest
windowed objects that \emph{are} residual, and Theorem~\ref{thm:packet} shows that this
includes arbitrary finite packet cardinalities.  Upgrading these finite packets to a
single infinite-horizon scrambled set requires diagonalizing across horizons in a
system-dependent way, reading unboundedly many maps to pin one system --- precisely the
``replacement for the backward induction consuming only finite-depth information per
step'' requested in \cite{companion}, now seen to also need to consume unboundedly many
\emph{times}.  We formulate the sharpened target.

\begin{problem}[Finite-horizon fusion]\label{prob:fusion}
Does there exist a comeager set on which a \emph{single} pair $(x,y)$ (depending on the
system) is DC1-scrambled at every tail with prescribed window visits at all times ---
equivalently, on which the finite-horizon witnesses of Theorem~\ref{thm:residual} can be
chosen coherently in $N$ so as to glue into an infinite-horizon scrambled pair?  By
Theorem~\ref{thm:obstruction} any such gluing cannot be certified by a bounded window of
times; it must be a genuinely non-finite-horizon (system-dependent, unbounded-window)
construction.
\end{problem}

A positive answer to Problem~\ref{prob:fusion} for the DC1 component, together with the
all-tail semifactor of Theorem~\ref{thm:dense}, would settle the DC1 half of
Problem~\ref{prob:companion}; the entropy and metric-mean-dimension halves would still
require their own non-finite-horizon arguments, by Theorem~\ref{thm:obstruction}.  We
intend to return to this elsewhere.

\subsection*{Acknowledgments}
The research presented in this paper was supported by the European Research Council
(ERC) under the European Union's Horizon Europe research and innovation programme
(grant agreement No.\ 101041711), by the Israel Science Foundation (grant number
2258/19) and by the Israel Science Foundation (ISF Grant 4101/25).

\subsection*{Declaration on generative AI}
During the preparation of this work the author used generative AI tools to accelerate
drafting and revision.  All mathematical claims, proofs, and bibliographic information
were subsequently reviewed and edited by the author, who takes full responsibility for
the content of the manuscript.

\begin{thebibliography}{9}

\bibitem{companion}
Y.~Shmalo,
\emph{Margin-stable schemes and the Baire category of full complexity in
non-autonomous dynamics on the Hilbert cube}, companion manuscript, 2026.

\bibitem{Shmalo2026}
Y.~Shmalo,
\emph{Strict cubical horseshoes and all-tail full complexity in non-autonomous
dynamics}, manuscript, submitted, 2026.

\bibitem{Miranda}
C.~Miranda,
\emph{Un'osservazione su un teorema di Brouwer},
Boll.\ Un.\ Mat.\ Ital.\ (2) \textbf{3} (1940--41), 5--7.

\bibitem{Kulpa}
W.~Kulpa,
\emph{The Poincar\'e--Miranda theorem},
Amer.\ Math.\ Monthly \textbf{104} (1997), no.~6, 545--550.

\bibitem{KolyadaSnoha}
S.~Kolyada and L.~Snoha,
\emph{Topological entropy of nonautonomous dynamical systems},
Random Comput.\ Dynam.\ \textbf{4} (1996), 205--233.

\bibitem{BalibreaProblems}
F.~Balibrea,
\emph{On problems of topological dynamics in non-autonomous discrete systems},
Appl.\ Math.\ Nonlinear Sci.\ \textbf{1} (2016), no.~2, 391--404.

\bibitem{Oxtoby}
J.~C.~Oxtoby,
\emph{The Banach--Mazur game and Banach category theorem},
in: Contributions to the Theory of Games, Vol.~III,
Ann.\ of Math.\ Stud.\ \textbf{39}, Princeton Univ.\ Press, 1957, 159--163.

\bibitem{LindenstraussWeiss}
E.~Lindenstrauss and B.~Weiss,
\emph{Mean topological dimension},
Israel J.\ Math.\ \textbf{115} (2000), 1--24.

\bibitem{SchweizerSmital}
B.~Schweizer and J.~Sm\'\i tal,
\emph{Measures of chaos and a spectral decomposition of dynamical systems on the
interval},
Trans.\ Amer.\ Math.\ Soc.\ \textbf{344} (1994), no.~2, 737--754.

\end{thebibliography}

\end{document}
